\section{Introduction}
\label{sec:introduction}

Air-quality policy in cities such as New Delhi needs to know which sources produce the
PM\(_{2.5}\) measured at the monitors, and published source apportionments for the same
Indian city disagree \citep{pant2012critical,karagulian2015contributions}.  Receptor
models estimate contributions from the chemical composition of particulate samples
\citep{watson2002receptor,paatero1994positive}, whereas we combine the hourly
PM\(_{2.5}\) and wind records of the New Delhi regulatory network
\citep{cpcb,cpcb_portal} with proxy emission maps for a few declared source groups
(brick kilns, industries, population-related activity, traffic) and a transport model
driven by the observed wind.

The readings then follow
a nonnegative linear model whose columns are the PM\(_{2.5}\) that one unit of a source
group's activity produces at every sensor and hour, plus a low-rank nuisance term for
regional background and smooth trends (Section~\ref{sec:problem_setup}).

A fit of this model with a small residual does not determine the contributions: a
strong distant source and a weaker nearby source on the same wind line produce nearly
proportional signals at the sensors, the nuisance term can absorb a source's signal,
and a sparse network may not see a source at all.  Very different splits then fit the
readings equally well, and a least-squares fit reports one regardless.  In our
controlled experiments, a fit without background removal reports group contributions
with errors of \(36\)--\(784\%\) (Section~\ref{sec:application}).

We therefore compute, before fitting, which groups of sources the network can
separate, and report each group's contribution and share only at that resolution.

\noindent\textbf{Contributions.}
\begin{itemize}[leftmargin=*,itemsep=0pt,topsep=0pt]
  \item We formulate PM\(_{2.5}\) attribution from a regulatory network as a nonnegative
  linear inverse problem (Section~\ref{sec:problem_setup}) and give a procedure,
  Identifiability-Aware Source Attribution (IASA), that fixes the reported grouping
  before fitting and reports group signals, shares, error bounds and intervals
  (Section~\ref{sec:method}).
  \item We show that a grouping of sources is identifiable if and only if the column
  spaces of its blocks form a direct sum, and that the finest such grouping is unique
  and computable in \(O(NJ^2)\) operations (Theorem~\ref{thm:groupings}).  We also show
  that the error of a group's fitted signal is at most the noise in the operator's
  column space divided by the group's separation from the other sources
  (Theorem~\ref{thm:group_error}), and we give a condition under which the set of
  reportable groupings is unchanged by an error in the wind-driven operator
  (Proposition~\ref{prop:grouping_stability}, Section~\ref{sec:identifiability_theory}).
  \item In controlled experiments with known contributions on the New Delhi domain,
  IASA flags an absorbed source and a pair of nearly proportional sources that no
  baseline flags, with group-signal errors of \(0.4\)--\(6.0\%\) (plain NNLS after
  background removal: \(0.4\)--\(6.2\%\), no flag)
  (Section~\ref{subsec:pollution_controlled}).  On four observed weeks from the
  32-sensor network the four source groups are separable in every week, but at the
  instrument noise level the error bound of the largest group's signal is below its
  fitted norm in only two of the four weeks
  (Section~\ref{subsec:observed_results}).
\end{itemize}
